\kchapter{Glossary and Learning Path}{k8s-logo}

\section{Glossary}

\begin{description}[style=nextline, leftmargin=1.2em, itemsep=0.35em, font=\color{k8sblue}\bfseries]
  \item[Admission controller] Code in the API server that can modify or reject
        requests after authentication and authorisation (e.g.\ Pod Security,
        ValidatingAdmissionPolicy).
  \item[Annotation] Non-identifying key/value metadata used by tools and humans.
  \item[API server] The control-plane component that serves the Kubernetes API;
        the only component that talks to etcd.
  \item[CNI] Container Network Interface: the plugin standard for Pod networking.
  \item[ConfigMap] An object holding non-secret configuration data.
  \item[Container runtime] Software that runs containers on a node (containerd, CRI-O).
  \item[Context] A kubeconfig entry combining a cluster, a user, and a default namespace.
  \item[Controller] A loop that watches objects and reconciles actual state
        toward desired state.
  \item[CRD] CustomResourceDefinition: adds a new resource type to the API.
  \item[CronJob] Creates Jobs on a schedule.
  \item[CSI] Container Storage Interface: the plugin standard for storage.
  \item[DaemonSet] Runs one Pod per (selected) node.
  \item[Deployment] Manages ReplicaSets to run and update stateless Pods.
  \item[Desired state] What you declared in an object's \kc{spec}.
  \item[EndpointSlice] The list of Pod IPs/ports currently backing a Service.
  \item[etcd] The consistent key--value store holding all cluster state.
  \item[Gateway API] The modern, role-oriented API for L4/L7 traffic routing
        (GatewayClass, Gateway, HTTPRoute\ldots).
  \item[GitOps] Operating model in which Git holds desired state and an in-cluster
        agent continuously applies it.
  \item[Helm] The Kubernetes package manager; packages are called charts.
  \item[HPA] HorizontalPodAutoscaler: scales the number of replicas on metrics.
  \item[Ingress] The original API for HTTP routing into the cluster.
  \item[Init container] A container that runs to completion before the app containers start.
  \item[Job] Runs Pods until a number of them complete successfully.
  \item[kubeconfig] The file \kc{kubectl} uses to find clusters and credentials.
  \item[kubelet] The node agent that runs Pods and reports their status.
  \item[kube-proxy] Programs node networking so Service virtual IPs reach Pods.
  \item[Kustomize] Template-free customisation of YAML through bases and overlays.
  \item[Label / selector] Identifying key/value pairs, and queries over them.
  \item[LimitRange] Per-namespace defaults and bounds for container resources.
  \item[Liveness probe] Health check whose failure restarts the container.
  \item[Namespace] A virtual partition of a cluster for names, RBAC, and quotas.
  \item[NetworkPolicy] Allow-list rules for Pod network traffic.
  \item[Node] A worker machine (VM or physical) in the cluster.
  \item[Operator] A controller that encodes operational expertise for specific software.
  \item[PDB] PodDisruptionBudget: limits voluntary disruptions of an app.
  \item[Pod] The smallest deployable unit: one or more containers sharing network and storage.
  \item[PV / PVC] PersistentVolume (actual storage) / PersistentVolumeClaim (a request for it).
  \item[QoS class] Guaranteed, Burstable, or BestEffort --- derived from requests and limits.
  \item[RBAC] Role-Based Access Control: Roles/ClusterRoles bound to subjects.
  \item[Readiness probe] Health check whose failure removes the Pod from Service endpoints.
  \item[ReplicaSet] Keeps a fixed number of identical Pods running.
  \item[Request / limit] Guaranteed / maximum amount of CPU or memory for a container.
  \item[ResourceQuota] A per-namespace cap on total resource consumption.
  \item[Scheduler] Assigns unscheduled Pods to nodes.
  \item[Secret] An object holding sensitive data (base64-encoded, not encrypted by default).
  \item[Service] A stable virtual IP and DNS name that load-balances to selected Pods.
  \item[ServiceAccount] An identity for processes running in Pods.
  \item[Sidecar] A helper container that runs alongside the main container for its whole life.
  \item[StatefulSet] Runs Pods with stable identities and per-Pod storage.
  \item[StorageClass] Describes a type of storage and enables dynamic provisioning.
  \item[Taint / toleration] Node-side repulsion and the Pod-side permission to ignore it.
\end{description}

\section{Your learning path}

\begin{center}
\resizebox{\linewidth}{!}{%
\begin{tikzpicture}[
  lvl/.style={draw=#1, fill=#1!8, rounded corners=6pt, line width=1pt,
              text width=4.5cm, minimum height=4.6cm, align=left, font=\footnotesize, text=k8snavy},
  hd/.style={fill=#1, text=white, rounded corners=4pt, font=\small\bfseries,
             minimum width=4.9cm, minimum height=0.8cm}]
  \node[lvl=k8sblue] (l1) at (0,0) {\begin{itemize}[leftmargin=1em]
        \item Containers \& Docker
        \item Pods, Deployments, Services
        \item ConfigMaps, Secrets
        \item kubectl fluency
        \item \textbf{Goal}: deploy an app with a Service
        \end{itemize}};
  \node[hd=k8sblue, anchor=south] at (l1.north) {\faSeedling\ \ Beginner};
  \node[lvl=k8sdeep] (l2) at (5.4,0) {\begin{itemize}[leftmargin=1em]
        \item Probes, resources, HPA
        \item StatefulSets, storage
        \item Gateway API, NetworkPolicy
        \item Helm, Kustomize
        \item \textbf{Cert}: CKAD
        \end{itemize}};
  \node[hd=k8sdeep, anchor=south] at (l2.north) {\faTree\ \ Intermediate};
  \node[lvl=k8snavy] (l3) at (10.8,0) {\begin{itemize}[leftmargin=1em]
        \item Cluster operations, upgrades
        \item RBAC, Pod Security, policy
        \item Observability, SLOs
        \item GitOps, CI/CD
        \item \textbf{Certs}: CKA, CKS
        \end{itemize}};
  \node[hd=k8snavy, anchor=south] at (l3.north) {\faMountain\ \ Advanced};
  \node[lvl=k8sorange] (l4) at (16.2,0) {\begin{itemize}[leftmargin=1em]
        \item Operators \& controllers
        \item Multi-cluster, service mesh
        \item Platform engineering
        \item Cost \& capacity
        \item \textbf{Path}: Kubestronaut
        \end{itemize}};
  \node[hd=k8sorange, anchor=south] at (l4.north) {\faRocket\ \ Expert};
\end{tikzpicture}}
\end{center}

\section{Certifications}

\begin{center}
\ktablestart
\begin{tabularx}{\linewidth}{P{2.0cm} L L}
\kheader \kth{Exam} & \kth{Focus} & \kth{Good preparation} \\
\textbf{KCNA} & Cloud-native fundamentals (multiple choice). & Parts I--III of this book. \\
\textbf{CKAD} & Building and running applications (hands-on). & Parts II--IV, and the capstone. \\
\textbf{CKA} & Administering clusters (hands-on). & Parts I, V, VI; kubeadm practice. \\
\textbf{CKS} & Securing clusters and workloads (hands-on, requires CKA). & Chapter 16 in depth, Falco, supply chain. \\
\end{tabularx}
\end{center}

\section{Where to keep learning}

\begin{itemize}
  \item \textbf{Official documentation}: \url{https://kubernetes.io/docs/} --- the tasks and tutorials sections especially.
  \item \textbf{Kubernetes the Hard Way} (Kelsey Hightower): build a cluster by hand to understand every component.
  \item \textbf{Gateway API docs}: \url{https://gateway-api.sigs.k8s.io/}.
  \item \textbf{CNCF Landscape}: \url{https://landscape.cncf.io/} --- the map of the ecosystem.
  \item \textbf{KubeCon talks} on YouTube --- real-world experience from people running Kubernetes at scale.
  \item \textbf{Play}: break your kind cluster on purpose, then fix it. That is how experts are made.
\end{itemize}

\vfill
\begin{center}
\kimg[2.2cm]{k8s-logo}\\[0.8em]
{\Large\bfseries\color{k8sblue} Happy sailing, helmsman.}\\[0.3em]
{\small\color{k8sgray} You now know how to steer.}
\end{center}
